% V2主教材扫描PDF98--99思考题15.1--15.14，原图15.17--15.19待重绘。
\section*{第15章\quad 静电场中的电介质：思考题}
\paragraph{思考题 15.1}
地球相对无穷远的自电容约700微法，为什么实验室的小体积电容器能达到1000微法乃至更大？

\paragraph{思考题 15.2}
平行板电容公式C随板间距离d趋近于零而发散，为什么实际不能无限减小d来获得无限大电容？

\paragraph{思考题 15.3}
把额外电荷放在平行板电容器某一极板上，极板表面电荷如何重新分配？是否全部停留在原来朝向另一极板的一面？

\paragraph{思考题 15.4}
证明图15.17中有限平行板电容器电场不能在板边缘突然从常量跳到零，说明边缘场和电势连续性。

\paragraph{思考题 15.5}
若考虑真实平行板电容器的边缘场，实际电容与理想平行板近似相比通常更大还是更小？

\paragraph{思考题 15.6}
原图15.18中部分插入电介质板的电容器，在固定总电荷和固定外接电压两种条件下，介质倾向被吸入还是推出极板间？

\paragraph{思考题 15.7}
原图15.19中同一平行板周围增加各电位与电容器两极板相等的保护环，为什么可减弱边缘场造成的电容量测量偏差？

\paragraph{思考题 15.8}
两介质介面法向极化强度分别为P1和P2，求该界面束缚面电荷密度的边界表达式和符号约定。

\paragraph{思考题 15.9}
在给定自由电荷分布和介电边界的静电系统里，移动一块电介质，哪些量不变：D分布、E分布、闭合曲面D通量、闭合曲面E通量？

\paragraph{思考题 15.10}
由极性分子组成的液态介质，相对介电常数在升温时为什么通常下降？

\paragraph{思考题 15.11}
带正电的胶木棒为何能吸起原本电中性的纸屑？

\paragraph{思考题 15.12}
从真空电场高斯定律与束缚电荷密度关系出发，推导电位移D的高斯定律，并说明它统计哪一类电荷。

\paragraph{思考题 15.13}
利用电容器能量或场能密度求理想长圆柱同轴电容器的单位长度电容，并注明边缘忽略。

\paragraph{思考题 15.14}
一根体电阻率\(\rho\) 的长直电阻导线，截面积A、稳态电流I，求导线表面外侧所需电场的大小与方向，以及其与电流做功的关系。
